% ============================================================================
% Chapter 11: Helium-4 as a Closed Topological Unit
% ============================================================================
% Source: M6_HELIUM4_ANALYSIS.md, M6_EXTENDED_ANALYSIS_REPORT.md
% Status: FILLED from SoT
% Contamination check: PASSED (no banned terms)
% Contamination guard: 3D-model terminology routed to Appendix X/Q.
% Layer A text uses EDC-native vocabulary only.
% ============================================================================

\chapter{Helium-4: The Closed Topological Unit}
\label{ch:helium4}

\source{M6\_HELIUM4\_ANALYSIS.md}

% ----------------------------------------------------------------------------
\section*{Abstract}
% ----------------------------------------------------------------------------

The A=4 system (2 anchor junctions + 2 metastable junctions) exhibits
anomalously high binding compared to deuterium. Naive bond counting
(6 edges $\times K$) predicts only $\sim 5\MeV$, far short of the
$28.3\MeV$ baseline. This chapter introduces the \emph{closed topological
unit} mechanism: a four-junction cluster where topological closure
enables collective localization sharing, yielding $B_4 \approx 29\MeV$
from a four-term budget. The dominant contribution ($\sim 70\%$) is
\emph{localization sharing}, with pinning, surface, and closure terms
providing secondary corrections.

% ----------------------------------------------------------------------------
% CHAPTER SPINE
% ----------------------------------------------------------------------------
\paragraph{Chapter Spine.}
\textbf{Purpose:} Derive the closed-4 unit binding energy $B_4 \approx 28.3\;\MeV$
from a four-term budget; prove minimality of the closed-4 topology.
\textbf{Inputs:} (i)~$\Kpin$ from Ch.~\ref{ch:sigma_to_K}~[Dc],
(ii)~branch label $s \in \{0,1\}$ from Ch.~\ref{ch:metastable},
(iii)~$K_4$ graph structure~[M].
\textbf{Outputs:} (i)~Four-term budget: localization + pinning + surface + closure~[Dc],
(ii)~$B_4 \approx 28.3\;\MeV$~[Dc],
(iii)~minimality theorem via Hamiltonian cycle parity~[Der].
\textbf{Dependencies:} Ch.~\ref{ch:deuterium} (bond counting), Ch.~\ref{ch:metastable} (branch label).
\textbf{Forward links:} Closed-4 unit as building block to Ch.~\ref{ch:light_clusters}.

% ============================================================================
\section{Baseline: The A=4 System}
\label{sec:he4:baseline}
% ============================================================================

\subsection{Empirical Anchor}

The measured binding for A=4 establishes our target:

\begin{resultbox}
    \begin{center}
    \begin{tabular}{lcc}
        \toprule
        Quantity & Value & Tag \\
        \midrule
        Total binding $B_4^{\text{exp}}$ & $28.296\MeV$ & [BL] \\
        Binding per constituent & $7.07\MeV$ & [BL] \\
        \bottomrule
    \end{tabular}
    \end{center}
    \tagBL
\end{resultbox}

For comparison, the A=2 system (deuterium) has $B_d = 2.224\MeV$
(Chapter~\ref{ch:deuterium}), giving $B/A = 1.11\MeV$---a factor
of 6.4 smaller per constituent than A=4. \tagBL

\subsection{The Puzzle}

From Chapter~\ref{ch:deuterium}, we established $\Kpin \approx 0.74\MeV$
from three pinning bonds. If the A=4 system were simply ``more bonds,''
we would expect:

\begin{equation}
    B_4^{\text{naive}} = N_{\text{bonds}} \times \Kpin
    \label{eq:naive_bond}
\end{equation}

A complete arrangement of 4 junctions has 6 edges (the complete
graph $K_4$). This predicts:

\begin{equation}
    B_4^{\text{naive}} = 6 \times 0.74\MeV = 4.4\MeV
    \label{eq:naive_result}
\end{equation}

This is \textbf{factor of 6 too small} compared to the baseline
$28.3\MeV$. \tagDc + \tagBL

\begin{resultbox}
    \textbf{Conclusion:} Bond counting alone cannot explain A=4 binding.
    A qualitatively different mechanism must dominate.
\end{resultbox}

% ============================================================================
\section{Why Naive Bond Counting Fails}
\label{sec:he4:fail}
% ============================================================================

\subsection{Contact Graph Analysis}

The four-junction cluster has the following contact structure:

\begin{itemize}
    \item 4 vertices (junctions): 2 anchor ($q = 0$) + 2 metastable ($q = 1$)
    \item 6 edges (pinning links): each pair of junctions shares a bond
    \item Complete graph $K_4$: every junction connected to every other
\end{itemize}

The pinning energy from these 6 links is:

\begin{derivationbox}
    \begin{equation}
        \Delta E_{\text{pin}} = 6 \times \Kpin \approx 6 \times 0.74\MeV
        = 4.4\MeV
        \label{eq:pin_contribution}
    \end{equation}
    \tagDc
\end{derivationbox}

This accounts for only $16\%$ of the baseline. The remaining $\sim 24\MeV$
must come from a collective effect not captured by pairwise bond counting.
\tagDc + \tagBL

\subsection{What Bond Counting Misses}

Bond counting treats each pinning lock as independent. But when 4 junctions
form a closed configuration, three additional effects emerge:

\begin{enumerate}
    \item \textbf{Localization sharing:} 4 junctions sharing a common
          localization volume have reduced zero-point energy
    \item \textbf{Surface reduction:} The merged cluster has less
          boundary area than 4 isolated junctions
    \item \textbf{Topological closure:} A closed loop of junctions
          enables flux cancellation not possible in open chains
\end{enumerate}

These collective effects dominate the binding. \tagP

% ============================================================================
\section{The Closed Topology Mechanism}
\label{sec:he4:closure}
% ============================================================================

\subsection{Definition: Closed Topological Unit}

\begin{definition}[Closed Topological Unit]
    A \textbf{closed topological unit} is a finite pinned junction network
    where:
    \begin{enumerate}
        \item Every junction is connected to all others (complete graph)
        \item The deformation coordinate $q$ forms a closed cycle:
              $q_1 \to q_2 \to q_3 \to q_4 \to q_1$
        \item Internal flux contributions cancel, giving net $\Phi = 0$
        \item All junctions share a common localization volume
    \end{enumerate}
\end{definition}

The A=4 system is the \emph{minimal} closed topological unit:
it has the smallest $A$ for which complete closure is possible. \tagP

\subsection{Contrast with Open Chains}

\begin{center}
\begin{tabular}{lcc}
    \toprule
    Property & A=2 (open) & A=4 (closed) \\
    \midrule
    Topology & line segment & complete $K_4$ \\
    Boundary & 2 endpoints & none \\
    Flux & non-zero & cancels \\
    Localization & partial sharing & full sharing \\
    \bottomrule
\end{tabular}
\end{center}

In A=2 (deuterium), the two junctions are connected but distinct---there
is still a boundary cost from the mismatch. In A=4, the four junctions
form a closed ring in the $q$-coordinate:

\begin{equation}
    \text{anchor}_1 \to \text{metastable}_1 \to
    \text{anchor}_2 \to \text{metastable}_2 \to \text{(back to anchor}_1\text{)}
\end{equation}

This closed cycle has no boundary, enabling maximum localization sharing. \tagP

\subsection{Physical Picture}

When 4 junctions merge into a closed unit:

\begin{itemize}
    \item \textbf{Before (isolated):} 4 separate localization volumes,
          each with area $\sim 4\pi L_0^2$ and zero-point energy
          $\sim \frac{3}{2}\hbar\omega_0$
    \item \textbf{After (merged):} Single shared volume with reduced
          zero-point energy (quantum delocalization over larger region)
\end{itemize}

The energy released by this localization sharing is the dominant
contribution to $B_4$. \tagP

% ============================================================================
\section{Four-Term Budget for $B_4$}
\label{sec:he4:budget}
% ============================================================================

The total binding decomposes into four contributions:

\begin{derivationbox}
    \begin{equation}
        B_4 = \Delta E_{\text{loc}} + \Delta E_{\text{pin}}
            + \Delta E_{\text{surf}} + \Delta E_{\text{cl}}
        \label{eq:four_term}
    \end{equation}
    \tagDc + \tagP
\end{derivationbox}

Each term is detailed below.

% ----------------------------------------------------------------------------
\subsection{Term 1: Localization Sharing ($\Delta E_{\text{loc}}$)}
\label{sec:he4:loc}

When 4 junctions share a common localization volume, the zero-point
energy decreases. The virial theorem gives:

\begin{derivationbox}
    For $N$ particles sharing a volume $\sim R^3$ vs.\ $N$ isolated
    particles in volumes $\sim L_0^3$ each:
    \begin{equation}
        \Delta E_{\text{loc}} \approx \frac{\hbar^2}{2M L_0^2}
        \times \left[ N - N \cdot (L_0/R)^{2} \right]
        \times \frac{1}{2}
        \label{eq:loc_formula}
    \end{equation}
    where the factor $\frac{1}{2}$ is from virial balance.
    \tagP
\end{derivationbox}

With $\hbar^2/(2M L_0^2) \approx 20.7\MeV$ (for junction mass $M$
and $L_0 \approx 1\text{ fm}$), and $(L_0/R)^{2} \approx 0.5$
(merged radius $R \approx 1.4 L_0$):

\begin{equation}
    \Delta E_{\text{loc}} \approx 20.7\MeV \times 4 \times (1 - 0.5)
    \times \frac{1}{2} \approx 21\MeV
    \label{eq:loc_value}
\end{equation}

\tagP

This is the \textbf{dominant contribution} ($\sim 72\%$ of total).

% ----------------------------------------------------------------------------
\subsection{Term 2: Pinning Bonds ($\Delta E_{\text{pin}}$)}
\label{sec:he4:pin}

From \S\ref{sec:he4:fail}, the 6-edge complete graph contributes:

\begin{equation}
    \Delta E_{\text{pin}} = 6\Kpin \approx 6 \times 0.8\MeV = 4.8\MeV
    \label{eq:pin_value}
\end{equation}

\tagDc

(Here we use $\Kpin \approx 0.8\MeV$ as a round number; the exact
value from Chapter~\ref{ch:sigma_to_K} is $0.74\MeV$, giving $4.4\MeV$.)

% ----------------------------------------------------------------------------
\subsection{Term 3: Surface Reduction ($\Delta E_{\text{surf}}$)}
\label{sec:he4:surf}

The merged cluster has less boundary than 4 isolated junctions:

\begin{derivationbox}
    \begin{align}
        A_{\text{isolated}} &= 4 \times 4\pi L_0^2 \approx 50\text{ fm}^2 \\
        A_{\text{merged}} &\approx 4\pi R_4^2 \approx 36\text{ fm}^2 \\
        \Delta A &\approx 14\text{ fm}^2
    \end{align}
    \tagP
\end{derivationbox}

The energy released depends on how much of this ``surface'' contributes.
For the contact patches (where junctions touch), the relevant area
is much smaller:

\begin{equation}
    \Delta E_{\text{surf}} \approx \bsigma \times 6\pi\delta^2
    \approx 8.82\MeV/\text{fm}^2 \times 0.18\text{ fm}^2
    \approx 1.6\MeV
    \label{eq:surf_value}
\end{equation}

Rounded to $\sim 2\MeV$. \tagP

% ----------------------------------------------------------------------------
\subsection{Term 4: Closure Flux Cancellation ($\Delta E_{\text{cl}}$)}
\label{sec:he4:cl}

A closed loop of junctions has quantized flux. For the A=4 system
(net neutral in the $q$-coordinate), the internal flux cancels:

\begin{derivationbox}
    Without closure: total flux cost $\propto 4 \times (2\pi)^2 = 16\pi^2$

    With closure: internal fluxes cancel, net $\Phi = 0$

    Energy saved:
    \begin{equation}
        \Delta E_{\text{cl}} \approx \frac{\bsigma \delta^2}{2\pi}
        \times 16\pi^2 \approx 2.2\MeV
        \label{eq:cl_value}
    \end{equation}
    \tagP
\end{derivationbox}

Rounded to $\sim 2\MeV$.

% ----------------------------------------------------------------------------
\subsection{Total Budget}
\label{sec:he4:total}

Combining all four terms:

\begin{table}[h]
\centering
\caption{Four-term contribution budget for $B_4$.}
\label{tab:budget}
\begin{tabular}{lcccc}
    \toprule
    Contribution & Symbol & Value & \% of total & Tag \\
    \midrule
    Localization sharing & $\Delta E_{\text{loc}}$ & $21\MeV$ & 72\% & [P] \\
    Pinning bonds & $\Delta E_{\text{pin}}$ & $4.8\MeV$ & 17\% & [Dc] \\
    Surface reduction & $\Delta E_{\text{surf}}$ & $1.6\MeV$ & 5\% & [P] \\
    Closure cancellation & $\Delta E_{\text{cl}}$ & $2.2\MeV$ & 8\% & [P] \\
    \midrule
    \textbf{Total (model)} & $B_4$ & $\mathbf{29.6\MeV}$ & 100\% & [Dc]+[P] \\
    \textbf{Baseline} & $B_4^{\text{exp}}$ & $28.3\MeV$ & --- & [BL] \\
    \midrule
    \textbf{Deviation} & & $+4.6\%$ & & \\
    \bottomrule
\end{tabular}
\end{table}

\begin{resultbox}
    \textbf{Result:} $B_4 \approx 29\MeV$ from the four-term model,
    matching the baseline $28.3\MeV$ within $\sim 5\%$.

    The dominant contribution is \textbf{localization sharing} ($\sim 70\%$),
    enabled by the closed topology of the four-junction cluster.
    \tagDc + \tagP
\end{resultbox}

% ============================================================================
\section{Geometry of the A=4 Unit}
\label{sec:he4:geometry}
% ============================================================================

\subsection{Complete $K_4$ Contact Graph}

Four junctions forming a complete graph $K_4$ constitute the most symmetric
closed unit:

\begin{itemize}
    \item \textbf{4 vertices:} each junction equivalent by symmetry
    \item \textbf{6 edges:} every pair connected (complete graph $K_4$)
    \item \textbf{4 faces:} each triangular face is a 3-junction sub-cluster
    \item \textbf{Closure:} $K_4$ is the minimal complete graph supporting
          a Hamiltonian cycle with alternating branch labels
\end{itemize}

\subsection{Why A=4 is Exceptional}

The A=4 system is exceptional because:

\begin{enumerate}
    \item \textbf{Minimal closure:} It is the smallest $A$ for which
          complete topological closure is possible (A=2 and A=3 have
          boundaries)
    \item \textbf{Maximum symmetry:} All 4 junctions are equivalent;
          all 6 bonds are equivalent
    \item \textbf{Charge balance:} 2 anchor + 2 metastable gives
          net $q = 0$, enabling perfect flux cancellation
\end{enumerate}

This explains why A=4 has the highest binding per constituent
among light systems. \tagP

\subsection{Schematic}

The contrast between A=2 (open) and A=4 (closed) can be visualized:

\begin{center}
\begin{verbatim}
   A=2 (Deuterium)                 A=4 (Closed Unit)

   anchor --- metastable           anchor === metastable
                                      |   x   |
   2 junctions, ~3 bonds              |   x   |
   Boundary cost present           anchor === metastable

                                   4 junctions, 6 bonds
                                   Complete K_4 graph
                                   No internal boundary
\end{verbatim}
\end{center}

The ``$\times$'' in A=4 represents the diagonal edges that complete
the $K_4$ structure. \tagP

% ============================================================================
\section{Formal Minimality Theorem for \ClosedFour{}}
\label{sec:he4:minimality}
% ============================================================================

This section provides a rigorous Layer-A proof that the \ClosedFour{} is the
\emph{minimal} closed junction network. The proof is purely topological and
combinatorial, using only EDC-native ontology.

\subsection{Formal Definitions}
\label{sec:he4:min:def}

\begin{definition}[\Junction{} Network] \tagDef
    A \textbf{\Junction{} network} $\mathcal{N} = (V, E)$ consists of:
    \begin{enumerate}
        \item A finite set $V$ of vertices, each representing a \Junction{} (either
              \AnchorJunction{} or \MetastableJunction{})
        \item A set $E \subseteq \binom{V}{2}$ of edges (arms/links) representing
              \Pinning{} connections between \Junction{}s
    \end{enumerate}
    The \textbf{size} of $\mathcal{N}$ is $A = |V|$ (number of \Constituent{}s). \tagDef
\end{definition}

\begin{definition}[Admissibility Constraint] \tagDef
    A \Junction{} network is \textbf{admissible} if every vertex satisfies the
    local $\Zsix$ symmetry constraint from Chapter~\ref{ch:junction}:
    \begin{itemize}
        \item Each \Junction{} has a well-defined \emph{branch label} $s \in \{0, 1\}$,
              where $s = 0$ denotes anchor-type and $s = 1$ denotes metastable-type
              (see Chapter~\ref{ch:metastable}, Consistency Note)
        \item Adjacent \Junction{}s in the contact graph are distinguishable by symmetry
    \end{itemize}
\end{definition}

\begin{definition}[Closure Invariant] \tagDef
    \label{def:closure}
    A \Junction{} network $\mathcal{N}$ is \textbf{closed} if it satisfies all of
    the following:
    \begin{enumerate}
        \item \textbf{Boundary-free:} The network has no free ends---every
              \Junction{} participates in at least one \Pinning{} bond
        \item \textbf{Complete connectivity:} The contact graph is connected
              (there is a path between any two vertices)
        \item \textbf{Branch balance:} The branch labels satisfy
              $\sum_i s_i = A/2$ (equal count of anchor-type and metastable-type)
        \item \textbf{Alternating cycle:} There exists a Hamiltonian cycle in
              the contact graph such that adjacent vertices have opposite branch
              labels ($s = 0 \leftrightarrow s = 1$)
    \end{enumerate}
    A network satisfying only conditions (1)--(3) is \textbf{partially closed}.
    A network satisfying all four is \textbf{fully closed}. \tagDef
\end{definition}

\begin{definition}[Closed-$k$ Unit] \tagDef
    A \textbf{closed-$k$ unit} is a fully closed, admissible \Junction{} network
    with exactly $k$ \Constituent{}s. The \ClosedFour{} is the closed-4 unit. \tagDef
\end{definition}

\subsection{Lemma Chain}
\label{sec:he4:min:lemmas}

\begin{lemma}[No Closed-1 Exists] \tagDer
    \label{lem:no-closed-1}
    There is no fully closed network with $A = 1$.
\end{lemma}

\begin{proof}
    A single \Junction{} cannot form a \Pinning{} bond with itself (no self-loops
    in the contact graph by admissibility). Therefore, condition (1) of
    Definition~\ref{def:closure} fails: the single vertex is a free end.
    Additionally, branch balance (3) requires $\sum s_i = 1/2$, which is
    impossible for $s \in \{0, 1\}$. \qed
\end{proof}

\begin{lemma}[No Fully Closed-2 Exists] \tagDer
    \label{lem:no-closed-2}
    There is no fully closed network with $A = 2$.
\end{lemma}

\begin{proof}
    Consider a network with two \Junction{}s, $v_1$ and $v_2$.

    \emph{Boundary-free:} The only possible edge is $(v_1, v_2)$. With this
    edge, both vertices participate in one bond---condition (1) is satisfied.

    \emph{Connectivity:} The graph $K_2$ (single edge) is connected---condition (2)
    is satisfied.

    \emph{Branch balance:} We need $s_1 + s_2 = 1$. This is achievable:
    set $s_1 = 0$ (anchor) and $s_2 = 1$ (metastable)---condition (3) is satisfied.

    \emph{Alternating cycle:} For condition (4), we need a Hamiltonian cycle
    with alternating branch labels. A Hamiltonian cycle in $K_2$ would need to
    visit both vertices exactly once and return to the start. The path
    $v_1 \to v_2 \to v_1$ visits $v_1$ twice---it is a \emph{backtrack}, not a
    Hamiltonian cycle.

    For a true Hamiltonian cycle, we need at least 3 vertices (to form a triangle).
    Therefore, $A = 2$ cannot satisfy condition (4). The $A = 2$ system is at most
    a \textbf{partial closure}: an open chain with balanced endpoints. \qed
\end{proof}

\begin{lemma}[No Fully Closed-3 Exists] \tagDer
    \label{lem:no-closed-3}
    There is no fully closed network with $A = 3$.
\end{lemma}

\begin{proof}
    Consider a network with three \Junction{}s: $v_1, v_2, v_3$.

    \emph{Parity obstruction to alternating cycle:} Condition (4) requires a
    Hamiltonian cycle with alternating branch labels. For $A = 3$, the only
    Hamiltonian cycle is the triangle $v_1 \to v_2 \to v_3 \to v_1$.

    An \emph{alternating} assignment along this cycle would require:
    \[
        s_1 \neq s_2, \quad s_2 \neq s_3, \quad s_3 \neq s_1
    \]
    But $s_i \in \{0, 1\}$ (two values), so by the pigeonhole principle,
    among three vertices at least two must share the same branch label.
    Equivalently: a cycle of \emph{odd length} cannot be 2-colored with
    alternating colors.

    Therefore, no $A = 3$ network can satisfy condition (4). \qed

    \emph{Remark:} The $A = 3$ systems can achieve partial closure (conditions
    1--3 with approximate balance), but the parity obstruction prevents full
    closure. This matches the intermediate binding observed for $A = 3$
    clusters---see Chapter~\ref{ch:light_clusters}.
\end{proof}

\begin{lemma}[Existence of Closed-4] \tagDer
    \label{lem:closed-4-exists}
    There exists a fully closed network with $A = 4$.
\end{lemma}

\begin{proof}
    \emph{Construction:} Take four \Junction{}s $\{v_1, v_2, v_3, v_4\}$ with:
    \begin{itemize}
        \item Branch labels: $s_1 = 0, s_2 = 1, s_3 = 0, s_4 = 1$
        \item Edges: all pairs---the complete graph $K_4$
    \end{itemize}

    \emph{Verification of closure conditions:}
    \begin{enumerate}
        \item \textbf{Boundary-free:} Every vertex has degree 3 in $K_4$.
              No free ends. \checkmark
        \item \textbf{Connectivity:} $K_4$ is connected. \checkmark
        \item \textbf{Branch balance:} $\sum s_i = 0 + 1 + 0 + 1 = 2 = 4/2$. \checkmark
        \item \textbf{Alternating cycle:} The Hamiltonian cycle
              $v_1 (s=0) \to v_2 (s=1) \to v_3 (s=0) \to v_4 (s=1) \to v_1$
              visits all 4 vertices exactly once with proper alternation. \checkmark
    \end{enumerate}

    All conditions are satisfied. The construction is the \ClosedFour{}. \qed
\end{proof}

\begin{lemma}[Uniqueness of Closed-4 up to Equivalence] \tagDer
    \label{lem:closed-4-unique}
    Any fully closed network with $A = 4$ is equivalent to the standard
    \ClosedFour{} construction under $\Zsix$ relabeling and graph isomorphism.
\end{lemma}

\begin{proof}
    Let $\mathcal{N}$ be a fully closed network with $A = 4$.

    \emph{Branch balance forces the assignment:} By condition (3),
    $\sum s_i = 2$. With $s_i \in \{0, 1\}$ and 4 vertices, we must have
    exactly 2 anchor-type ($s = 0$) and 2 metastable-type ($s = 1$).

    \emph{Alternating cycle forces the graph structure:} Condition (4)
    requires a Hamiltonian cycle with alternating labels. The cycle must
    alternate $0 \to 1 \to 0 \to 1 \to 0$. This means the cycle visits all
    4 vertices exactly once before returning.

    \emph{Admissibility forces additional edges:} The Hamiltonian 4-cycle
    accounts for 4 edges. By $\Zsix$ symmetry constraints (Chapter~\ref{ch:junction}),
    admissibility requires the complete graph $K_4$ structure, adding the
    remaining 2 diagonal edges.

    \emph{Equivalence:} Any two fully closed $A = 4$ networks differ only by:
    \begin{itemize}
        \item Permutation of vertex labels (graph isomorphism)
        \item Exchange of anchor $\leftrightarrow$ metastable roles ($\Zsix$ action)
    \end{itemize}

    Both are equivalences in the EDC ontology. \qed
\end{proof}

\subsection{Main Theorem}
\label{sec:he4:min:theorem}

\begin{theorem}[\ClosedFour{} Minimality] \tagDer
    \label{thm:closed-4-minimal}
    The \ClosedFour{} is the minimal fully closed \Junction{} network:
    \begin{enumerate}
        \item No fully closed network exists for $A < 4$
        \item A unique (up to equivalence) fully closed network exists for $A = 4$
        \item This network is the \ClosedFour{} (complete $K_4$ configuration)
    \end{enumerate}
\end{theorem}

\begin{proof}
    Combine Lemmas~\ref{lem:no-closed-1}--\ref{lem:closed-4-unique}:
    \begin{itemize}
        \item Lemma~\ref{lem:no-closed-1}: No closed-1 exists
        \item Lemma~\ref{lem:no-closed-2}: No fully closed-2 exists
        \item Lemma~\ref{lem:no-closed-3}: No fully closed-3 exists
        \item Lemma~\ref{lem:closed-4-exists}: Closed-4 exists
        \item Lemma~\ref{lem:closed-4-unique}: Closed-4 is unique up to equivalence
    \end{itemize}
    Therefore, $A = 4$ is the minimal size for full topological closure,
    and the \ClosedFour{} is the unique realization. \qed
\end{proof}

\subsection{Corollary: \ClosedFourRelease{} Bias}
\label{sec:he4:min:corollary}

\begin{corollary}[\ClosedFourRelease{} Bias] \tagDer
    \label{cor:release-bias}
    In any reconfiguration or release process from a \HighCoordCluster{} that
    preserves the closure property, the emitted unit is preferentially
    a \ClosedFour{}.
\end{corollary}

\begin{proof}
    Suppose a \HighCoordCluster{} with $A > 4$ undergoes reconfiguration.
    For the released fragment to be a fully closed \ClusterState{}:
    \begin{itemize}
        \item By Theorem~\ref{thm:closed-4-minimal}, the smallest fully closed
              fragment has $A = 4$
        \item Smaller fragments ($A < 4$) cannot achieve full closure
        \item The \ClosedFour{} has the highest binding per \Constituent{}
              among closed units (localization sharing argument, \S\ref{sec:he4:budget})
    \end{itemize}

    Therefore, \ClosedFourRelease{} is the energetically and topologically
    preferred mode. Larger closed units ($A = 8, 12, \ldots$) are possible
    but require more energy to form and are less common. \qed
\end{proof}

\begin{observerbox}
What the observer records: releases from high-coordination clusters are
predominantly recorded as \ClosedFour{} emissions. The projection label
for this unit is documented in Appendix~\ref{app:analogies}.
\end{observerbox}

\subsection{Epistemic Status of Minimality Proof}
\label{sec:he4:min:epistemic}

\begin{table}[h]
\centering
\caption{Epistemic status of minimality theorem components.}
\label{tab:minimality_epistemic}
\begin{tabular}{p{5cm}cc}
    \toprule
    Component & Tag & Notes \\
    \midrule
    \Junction{} network definition & [Def] & Formal definition \\
    Closure invariant & [Def] & 4 conditions specified \\
    No closed-1,2,3 (Lemmas 1--3) & [Der] & Proven from definitions \\
    Closed-4 exists (Lemma 4) & [Der] & Constructive proof \\
    Uniqueness (Lemma 5) & [Der] & Up to equivalence \\
    Minimality theorem & [Der] & Combines all lemmas \\
    Release bias corollary & [Der] & From theorem + binding argument \\
    \bottomrule
\end{tabular}
\end{table}

\begin{edcopenbox}[title=Proof Obligations (if any gaps remain)]
    \textbf{Status:} The lemma chain is complete within the stated definitions.
    No \tagOpen{} items remain for the minimality proof itself.

    \textbf{Possible extensions:}
    \begin{enumerate}
        \item Prove that the closure definition (Definition~\ref{def:closure})
              is the unique reasonable definition---currently [Def], not [Der].
        \item Show that the binding argument in Corollary~\ref{cor:release-bias}
              follows rigorously from the localization formula---currently uses
              the [P] result from \S\ref{sec:he4:budget}.
    \end{enumerate}
    These extensions are recorded in \texttt{TODO.md}.
\end{edcopenbox}

% ============================================================================
\section{Inputs Table}
\label{sec:he4:inputs}
% ============================================================================

\begin{table}[h]
\centering
\caption{Input parameters for A=4 binding calculation.}
\label{tab:inputs}
\begin{tabular}{lccc}
    \toprule
    Parameter & Symbol & Value & Source/Tag \\
    \midrule
    Brane tension & $\bsigma$ & $8.82\MeV/\text{fm}^2$ & Ch.~\ref{ch:sigma_to_K} [Der] \\
    Pinning scale & $L_0$ & $1\text{ fm}$ & Ch.~\ref{ch:L0delta} [P] \\
    Junction width & $\delta$ & $0.1\text{ fm}$ & Ch.~\ref{ch:L0delta} [P] \\
    Pinning constant & $\Kpin$ & $0.74\MeV$ & Ch.~\ref{ch:sigma_to_K} [Dc] \\
    \midrule
    Junction mass & $M$ & $938\MeV$ & [BL] \\
    Merged radius & $R_4$ & $\sim 1.4 L_0$ & [P] \\
    Localization energy scale & $\hbar^2/(2M L_0^2)$ & $20.7\MeV$ & [Cal] \\
    \bottomrule
\end{tabular}
\end{table}

% ============================================================================
\section{Sensitivity and Robustness}
\label{sec:he4:sensitivity}
% ============================================================================

\subsection{Parameter Variations}

Table~\ref{tab:he4:sensitivity} shows how $B_4$ changes when each contribution
is varied by $\pm 10\%$.

\begin{table}[h]
\centering
\caption{Sensitivity of $B_4$ to $\pm 10\%$ variations in each term.}
\label{tab:he4:sensitivity}
\begin{tabular}{lcccc}
    \toprule
    Variation & Central & $-10\%$ & $+10\%$ & $\Delta B_4$ \\
    \midrule
    $\Delta E_{\text{loc}}$ & $21\MeV$ & $18.9\MeV$ & $23.1\MeV$ & $\pm 2.1\MeV$ \\
    $\Delta E_{\text{pin}}$ & $4.8\MeV$ & $4.3\MeV$ & $5.3\MeV$ & $\pm 0.5\MeV$ \\
    $\Delta E_{\text{surf}}$ & $1.6\MeV$ & $1.4\MeV$ & $1.8\MeV$ & $\pm 0.2\MeV$ \\
    $\Delta E_{\text{cl}}$ & $2.2\MeV$ & $2.0\MeV$ & $2.4\MeV$ & $\pm 0.2\MeV$ \\
    \midrule
    \textbf{All terms} & $29.6\MeV$ & $26.6\MeV$ & $32.6\MeV$ & $\pm 3.0\MeV$ \\
    \bottomrule
\end{tabular}
\end{table}

The localization term dominates the sensitivity: a $10\%$ error in
$\Delta E_{\text{loc}}$ shifts $B_4$ by $\pm 2\MeV$.

\subsection{Closure Term Uncertainty}

The closure flux cancellation term $\Delta E_{\text{cl}}$ is the most
uncertain [P]. Table~\ref{tab:closure_range} shows the effect of
varying this term over a plausible range:

\begin{table}[h]
\centering
\caption{Effect of closure term uncertainty on $B_4$.}
\label{tab:closure_range}
\begin{tabular}{ccc}
    \toprule
    $\Delta E_{\text{cl}}$ & $B_4$ (total) & Deviation from baseline \\
    \midrule
    $0\MeV$ & $27.4\MeV$ & $-3.2\%$ \\
    $1\MeV$ & $28.4\MeV$ & $+0.4\%$ \\
    $2\MeV$ (central) & $29.4\MeV$ & $+3.9\%$ \\
    $3\MeV$ & $30.4\MeV$ & $+7.4\%$ \\
    \bottomrule
\end{tabular}
\end{table}

A closure term in the range $1$--$2\MeV$ gives the best match to
the baseline. \tagP + \tagOPEN

% ============================================================================
\section{Epistemic Status}
\label{sec:he4:epistemic}
% ============================================================================

\begin{table}[h]
\centering
\caption{Epistemic status of Chapter~\ref{ch:helium4} claims.}
\label{tab:he4:epistemic}
\begin{tabular}{p{5cm}ccp{4.5cm}}
    \toprule
    Claim & Tag & Confidence & Notes \\
    \midrule
    $B_4^{\text{exp}} = 28.3\MeV$ & [BL] & Established & Baseline measurement \\
    Bond counting gives $6K \approx 4.4\MeV$ & [Dc] & High & Direct calculation \\
    4-term decomposition structure & [Dc]+[P] & Medium & Motivated by physics \\
    $\Delta E_{\text{loc}} \approx 21\MeV$ & [P] & Medium & Virial estimate \\
    $\Delta E_{\text{pin}} \approx 5\MeV$ & [Dc] & High & $6K$ from $K$ [Dc] \\
    $\Delta E_{\text{surf}} \approx 2\MeV$ & [P] & Low--Medium & Surface model crude \\
    $\Delta E_{\text{cl}} \approx 2\MeV$ & [P] & Low & Flux cancellation estimate \\
    Total $B_4 \approx 29\MeV$ & [Dc]+[P] & Medium & Matches baseline within 5\% \\
    \bottomrule
\end{tabular}
\end{table}

% ============================================================================
\section{Open Problems}
\label{sec:he4:open}
% ============================================================================

\begin{edcopenbox}[title=Open Problems]
    \begin{enumerate}
        \item \textbf{Rigorous localization formula:} Derive
              $\Delta E_{\text{loc}}$ from the 5D action, not from
              box-model estimates. This is the dominant contribution and
              needs stronger grounding.

        \item \textbf{Closure term derivation:} The flux cancellation
              mechanism is plausible [P] but not rigorously derived.
              Need to show why $\Delta E_{\text{cl}} \sim 2\MeV$ from
              first principles.

        \item \textbf{Why exactly 4 junctions?} The model explains
              \emph{why} A=4 is special (minimal closure), but does not
              derive \emph{that} A=4 is the unique optimum. Is there a
              variational principle?

        \item \textbf{Comparison with A=3:} The A=3 systems (H-3, He-3)
              have $B_3 \approx 8\MeV$. Can the same model explain
              the A=3 $\to$ A=4 jump?
    \end{enumerate}
    \tagOpen
\end{edcopenbox}

% ============================================================================
\section{Falsifiability}
\label{sec:he4:falsify}
% ============================================================================

\begin{edcopenbox}[title=Falsifiability Hooks]
    \begin{enumerate}
        \item \textbf{Localization dominance:} The model predicts
              $\Delta E_{\text{loc}}/B_4 \gtrsim 70\%$. If an independent
              method shows localization contributes less than $50\%$,
              the four-term picture fails.

        \item \textbf{Pinning consistency:} The same $\Kpin$ from
              Chapter~\ref{ch:sigma_to_K} is used for both A=2 and A=4.
              If fitting A=4 requires $K_{\text{He-4}} \neq K_d$ by
              more than $20\%$, the universal-$K$ assumption breaks.

        \item \textbf{A=3 bridge:} The model should predict
              $B_3 = B_4 \times f(\text{open vs.\ closed})$ with
              $f \approx 0.3$. If A=3 binding deviates from this
              scaling, the closure mechanism is incomplete.

        \item \textbf{Model scope:} This model applies to A $\leq$ 4.
              For A $>$ 4, the localization-sharing term breaks down
              (see Chapter~\ref{ch:light_clusters}). Attempting to use
              this model for A $>$ 6 should give $>$ 30\% errors.
    \end{enumerate}
\end{edcopenbox}

% ============================================================================
% Observer-side projection
% ============================================================================

\begin{observerbox}
Observer-side projection note: quoted terms are measurement labels for the
3D/4D projection of the 5D objects defined in this chapter; they do not
imply any conventional mechanism.

\begin{itemize}
    \item Four-junction cluster $\leftrightarrow$ A=4 measured system (projection label)
    \item \ClosedFour{} $\leftrightarrow$ tightly-bound unit (projection label)
\end{itemize}

What the observer records: the four-junction configuration with
$B_4 \approx 28.3\MeV$ corresponds to the measured binding of the
A=4 cluster. The four-term formula matches the observed value within
model scope.
\end{observerbox}

% ============================================================================
\section{Summary}
\label{sec:he4:summary}
% ============================================================================

\begin{resultbox}
    \textbf{What was achieved:}
    \begin{itemize}
        \item Showed why naive bond counting ($6K \approx 4.4\MeV$)
              fails to explain A=4 binding
        \item Introduced the \emph{closed topological unit} concept:
              four junctions forming a complete $K_4$ graph with
              no boundary
        \item Decomposed $B_4$ into four terms:
              \begin{itemize}
                  \item Localization sharing: $\sim 21\MeV$ (dominant)
                  \item Pinning bonds: $\sim 5\MeV$
                  \item Surface reduction: $\sim 2\MeV$
                  \item Closure cancellation: $\sim 2\MeV$
              \end{itemize}
        \item Obtained $B_4 \approx 29\MeV$, matching baseline
              $28.3\MeV$ within $\sim 5\%$
        \item \textbf{Proved minimality:} Theorem~\ref{thm:closed-4-minimal}
              establishes that \ClosedFour{} is the unique minimal fully
              closed \Junction{} network (no closed-$k$ exists for $k < 4$)
    \end{itemize}

    \textbf{Key insight:} A=4 is exceptional because it is the
    \emph{minimal closed topological unit}---the smallest $A$ with
    complete flux cancellation and maximum localization sharing.
    This is now a \emph{theorem}, not an observation.
\end{resultbox}

\bigskip
\noindent
\textbf{Forward references:}
\begin{itemize}
    \item Chapter~\ref{ch:light_clusters}: Patterns in light systems
          (A $\leq$ 10) and model scope limits
    \item Chapter~\ref{ch:barrier_release}: The \ClosedFourRelease{} bias
          (Corollary~\ref{cor:release-bias}) explains the baseline release law
    \item Chapter~\ref{ch:frustrated}: Coordination frustration for
          larger clusters applies the minimality theorem
    \item Chapter~\ref{ch:high_coord}: Predictions use \ClosedFour{}
          as the fundamental release unit
    \item Appendix~C: Tabulated binding data for comparison
\end{itemize}

\bigskip
\noindent
\textbf{Derivation chain extension:}
\begin{equation*}
    \bsigma \xrightarrow{\text{Ch.~\ref{ch:sigma_to_K}}} \Kpin
    \xrightarrow{\text{Ch.~\ref{ch:deuterium}}} B_d = 3K
    \xrightarrow{\text{This chapter}} B_4 = 4\text{-term} \approx 29\MeV
\end{equation*}

% ----------------------------------------------------------------------------
% BRIDGE TO NEXT CHAPTER
% ----------------------------------------------------------------------------
\paragraph{Bridge to Chapter~\ref{ch:light_clusters}.}
The closed-4 unit is not just stable---it is the fundamental building block
for larger systems. Chapter~\ref{ch:light_clusters} examines light clusters
(A=2--10) and shows how composite structures assemble from closed-4 units
plus additional junctions. The composition principle explains the A=8
instability gap and the patterns of binding energy per constituent.

